\chapter{Production Roadmap}
\label{ch:roadmap}

This chapter documents the enhancements required to move the platform from its current Docker Compose deployment to a production banking environment.

\textbf{Container orchestration.} The system currently runs 20 Docker Compose services on a single host --- sufficient for development and demonstration, but production requires horizontal scaling, automatic failover, and resource-aware scheduling across nodes. Kubernetes would replace Compose: each streaming service becomes a Deployment with a configurable replica count, and the Prometheus metrics endpoints already exposed can serve as liveness and readiness probes with no code changes. The model volume would migrate to a Persistent Volume Claim, preserving the retrain-without-rebuild pattern of Chapter~\ref{ch:mlops}. RedPanda, Redis, and PostgreSQL would run as StatefulSets with persistent storage, or be replaced by managed equivalents depending on the bank's infrastructure policy.

\textbf{CI/CD.} Code changes are currently tested manually and deployed by rebuilding images. A production pipeline would add three stages: \emph{continuous integration} running unit tests for the three core modules on every commit, validating that feature computation produces the expected 30-dimensional output, and linting; \emph{image build} on merge to main, tagging images with the commit hash for traceability between deployed code and source control; and \emph{staged deployment} to a staging environment first, running the E2E test script there and promoting to production only after validation. The core extraction pattern simplifies this considerably --- the pure cores can be unit-tested without Redis, PostgreSQL, or RedPanda, covering most business logic with fast isolated tests.

\textbf{Secrets management.} Database credentials, API keys, and the Airflow REST password currently live in a \code{.env} file excluded from version control. This is adequate for a single developer but insufficient for production, where secrets must be rotated, audited, and never stored in plaintext on disk. A solution such as HashiCorp Vault or a cloud-native equivalent would provide encrypted storage, automatic rotation, and access logging, with Kubernetes Secrets injecting credentials at runtime and eliminating environment files.

\textbf{Load testing.} The system has been validated on 243{,}000 events replayed at accelerated speed but not under sustained production-equivalent load. A load testing phase would establish each service's throughput ceiling, identify the first bottleneck, and determine the point at which consumer lag begins to accumulate. The existing simulator can serve as the load generator with minimal modification --- adjusting the inter-event delay to simulate peak branch hours across 150 branches --- and the existing Prometheus instrumentation already captures the metrics needed to evaluate the results.

\textbf{Data retention.} The schema currently retains all transactions and alerts indefinitely. Regulatory requirements and storage constraints make a retention policy necessary: Tunisian banking regulations require records to be preserved for a minimum period, after which they may be archived or purged. A retention strategy would partition \code{transactions} and \code{alerts} by month using PostgreSQL's native partitioning, enabling efficient archival of older partitions to compressed cold storage while keeping recent data queryable at full speed. The nightly profile rebuild would then be scoped to the retention window rather than the full history, keeping computation time bounded as the dataset grows.
